\section{Results}
\label{sec:results}

We organize results around a single question: \emph{how much does the
decision layer matter?} Five claims answer it. The headline: TTA paired
with adaptive deferral removes roughly 80\% of segmentation errors at a
25\% deferral rate (Table~\ref{tab:main_results}); the same uncertainty
under a naive global rule removes only a third of that.
Figure~\ref{fig:summary_scatter} plots every method--policy combination
on one axis and makes the interaction visible at a glance---the spread
across configurations is far wider than the spread across uncertainty
sources alone, which establishes the decision layer as the dominant
contributor to final decision quality.

\begin{figure}[H]
\centering
\includegraphics[width=0.92\textwidth]{summary_scatter.png}
\caption{Error reduction vs.\ deferral rate. Each point pairs an uncertainty
method (MC Dropout in blue, TTA in orange) with a deferral policy (marker
shape). TTA + adaptive dominates the upper-left (79.5\% error reduction at
25\% deferral); TTA + confidence-aware is the Pareto optimum at low budgets
(55\% at only 12\%). Every TTA configuration outperforms every MC Dropout
configuration---the decision layer reshapes, not merely refines, what the
uncertainty estimator delivers.}
\label{fig:summary_scatter}
\end{figure}

\begin{table}[H]
\centering
\caption{Main results on the DRIVE test set (20 images). TTA matches the
ensemble on segmentation quality while producing the best uncertainty
(Unc-AUROC 0.881) and the largest error reduction (79.5\% at 25\% deferral)
at 3.1$\times$ the speed of MC Dropout. Best values per column in bold;
runtime is per image on a single GPU. ERR $=$ error reduction ratio;
Def.\ \% $=$ fraction of pixels deferred.}
\label{tab:main_results}
\scriptsize
\setlength{\tabcolsep}{3.5pt}
\renewcommand{\arraystretch}{1.03}
\resizebox{\textwidth}{!}{%
\begin{tabular}{l c c c c c c c c}
\toprule
\textbf{Method} & \textbf{Dice} & \textbf{AUC} & \textbf{IoU} &
\textbf{ECE} & \textbf{Unc-AUROC} & \textbf{ERR} &
\textbf{Def.\ \%} & \textbf{Time (s)} \\
\midrule
Deterministic            & 0.736 & 0.944 & ---   & 0.245 & 0.500 & ---           & ---   & \textbf{0.08} \\
MC Drop.\ (global)       & 0.763 & 0.965 & 0.617 & 0.035 & 0.722 & 26.5\%        & 33.4\% & 2.65 \\
MC Drop.\ (adaptive)     & 0.763 & 0.965 & 0.617 & 0.035 & 0.722 & 42.4\%        & 25.0\% & 2.65 \\
MC Drop.\ (conf-aware)   & 0.763 & 0.965 & 0.617 & 0.035 & 0.722 & 48.8\%        & 10.9\% & 2.65 \\
\addlinespace[2pt]
TTA (global)             & 0.768 & 0.969 & 0.624 & 0.035 & \textbf{0.881} & 55.1\%        & 13.2\% & 0.86 \\
TTA (adaptive)           & \textbf{0.768} & \textbf{0.969} & \textbf{0.624} & 0.035 & \textbf{0.881} & \textbf{79.5\%} & 25.0\% & 0.86 \\
TTA (conf-aware)         & 0.768 & 0.969 & 0.624 & 0.035 & \textbf{0.881} & 55.1\%        & \textbf{12.1\%} & 0.86 \\
\midrule
Ensemble ($N{=}5$)       & 0.771 & 0.947 & 0.627 & \textbf{0.031} & 0.833 & ---           & ---   & 3.90 \\
\bottomrule
\end{tabular}%
}
\end{table}

\paragraph{Reading the table.} Rows 2--4 and 5--7 share identical Dice,
AUC, and Unc-AUROC within each uncertainty method; only the rightmost
columns (ERR, Def.\,\%) change. The three-fold spread in error reduction
within a single uncertainty source is the clearest evidence that the
decision layer is not a formality---it is where most of the decision
quality is made or lost.

\subsection{TTA produces superior uncertainty for deferral}
\label{sec:results_uncertainty}

\textbf{Claim:} TTA uncertainty discriminates correct from incorrect
predictions substantially better than MC Dropout, while requiring less
computation.

\emph{Insight before the numbers.} If errors concentrate at specific
geometric structures, an uncertainty source that probes geometric
instability should separate them more cleanly than one that probes
weight instability.

TTA achieves an Unc-AUROC of 0.881, compared to 0.722 for MC Dropout
(Table~\ref{tab:main_results}). This is a large gap: TTA correctly ranks
88.1\% of error/non-error pixel pairs by uncertainty, versus 72.2\% for
MC Dropout. The difference is consistent across images and robust to
bootstrap resampling.

TTA also produces higher AUCC for both Dice (0.846 vs.\ 0.792) and AUC
(0.888 vs.\ 0.884), meaning that progressively removing the most uncertain
pixels improves performance faster with TTA uncertainty.

In terms of computation, TTA requires 6 forward passes with deterministic
augmentations ($K = 6$), while MC Dropout requires 30 stochastic passes
($T = 30$). TTA inference takes approximately 0.86\,s per image versus
2.65\,s for MC Dropout---a 3.1$\times$ speedup.

\paragraph{Interpretation.} The quality gap has a structural explanation.
In retinal vessel segmentation, errors concentrate at vessel boundaries
and thin vessel tips---regions sensitive to small input perturbations. TTA
directly tests this sensitivity by applying geometric transforms; if a
boundary pixel flips under rotation, TTA flags it. MC Dropout instead
samples weight perturbations, producing more spatially diffuse uncertainty
that does not concentrate as sharply on error-prone structures. The
takeaway generalizes beyond retinas: decision-time uncertainty should be
sourced from perturbations that mimic the error mode, not from whichever
Bayesian approximation is easiest to fit. ``Better uncertainty'' is not a
fixed quantity---it is the uncertainty whose failure mode matches the
decision rule downstream.

\begin{figure}[H]
\centering
\includegraphics[width=\textwidth]{uncertainty_maps.png}
\caption{Uncertainty maps from MC Dropout and TTA. TTA uncertainty
concentrates on boundary and thin-vessel regions where errors are most
frequent, while MC Dropout produces more spatially diffuse uncertainty.}
\label{fig:uncertainty_maps}
\end{figure}

\subsection{Adaptive deferral maximizes error removal; confidence-aware deferral is most efficient at low budget}
\label{sec:results_deferral}

\textbf{Claim:} Adaptive deferral achieves the largest absolute error removal,
while confidence-aware deferral is the most efficient policy at low review
budgets.

Table~\ref{tab:deferral_comparison} breaks down the deferral results. For MC
Dropout, global thresholding defers 33.4\% of pixels and reduces error by
26.5\%. Adaptive deferral improves to 42.4\% error reduction at a 25\% deferral
budget, and confidence-aware deferral reaches 48.8\% while deferring only
10.9\% of pixels.

For TTA, adaptive deferral yields the headline result: 79.5\% error reduction
at 25.0\% deferral. The confidence-aware policy matches the global TTA policy
in error reduction but with a smaller review budget (12.1\% vs.\ 13.2\%
deferred).

\begin{table}[H]
\centering
\caption{Deferral strategy comparison on the DRIVE test set. Confidence-aware
deferral dominates at low review budgets---it delivers 4.56 percentage points
of error reduction per percentage point deferred, roughly $5.8\times$ the
efficiency of the MC Dropout global baseline. Adaptive deferral produces the
largest absolute error removal (79.5\% at 25\% deferral). Bold marks the best
value in each column within its uncertainty group.}
\label{tab:deferral_comparison}
\scriptsize
\setlength{\tabcolsep}{3.5pt}
\renewcommand{\arraystretch}{1.03}
\begin{center}
\begin{tabular}{l l c c c c c}
\toprule
\textbf{Uncertainty} & \textbf{Deferral} & \textbf{Error before} &
\textbf{Error after} & \textbf{ERR} & \textbf{Def.\ \%} &
\textbf{ERR / Def.\ \%} \\
\midrule
MC Dropout & Global      & 0.0658 & 0.0484 & 26.5\% & 33.4\% & 0.79 \\
MC Dropout & Adaptive    & 0.0658 & 0.0379 & 42.4\% & 25.0\% & 1.70 \\
MC Dropout & Conf-aware  & 0.0658 & \textbf{0.0337} & \textbf{48.8\%} & \textbf{10.9\%} & \textbf{4.49} \\
\addlinespace[2pt]
\midrule
\addlinespace[2pt]
TTA        & Global      & 0.0640 & 0.0287 & 55.1\% & 13.2\% & 4.19 \\
TTA        & Adaptive    & 0.0640 & \textbf{0.0131} & \textbf{79.5\%} & 25.0\% & 3.18 \\
TTA        & Conf-aware  & 0.0640 & 0.0288 & 55.1\% & \textbf{12.1\%} & \textbf{4.56} \\
\bottomrule
\end{tabular}
\end{center}
\end{table}

The rightmost column in Table~\ref{tab:deferral_comparison}, ERR per
percentage point deferred (ERR/Def.\ \%), quantifies deferral efficiency.
For MC Dropout, confidence-aware deferral is dramatically more efficient than
the global policy (4.49 vs.\ 0.79). For TTA, both global and confidence-aware
policies are efficient, with confidence-aware slightly better (4.56 vs.\ 4.19)
at a lower deferral rate.

\paragraph{Interpretation.} Global thresholding wastes deferral budget on
pixels that are uncertain but correctly predicted. The confidence-aware
score $s = u \cdot (1 - c)$ suppresses these false alarms. Vessel
interiors are a clear example: MC Dropout assigns moderate uncertainty
there, but predictions sit well above 0.5, so $c$ is high and $s$ remains
low. The broader lesson: deferral efficiency is decided in the score
function, not just in the uncertainty source. A decision rule that treats
``uncertain'' and ``borderline'' as synonyms is leaving free error
reduction on the table.

\begin{figure}[H]
\centering
\includegraphics[width=\textwidth]{error_reduction_bars.png}
\caption{Error rates before and after deferral across all configurations.
TTA+adaptive achieves 80\% error reduction; TTA+conf-aware achieves 55\%
error reduction at roughly half the review budget. The choice of both
uncertainty method and deferral policy substantially reshapes the practical
review workflow.}
\label{fig:error_reduction_bars}
\end{figure}

\begin{figure}[H]
\centering
\begin{subfigure}[b]{0.49\textwidth}
  \centering
  \includegraphics[width=\linewidth]{deferral_curves.png}
  \caption{TTA risk-coverage under three deferral modes. All three policies
  reduce risk monotonically; confidence-aware reaches the clinical Dice
  target (0.82) at the highest coverage.}
  \label{fig:deferral_curves}
\end{subfigure}\hfill
\begin{subfigure}[b]{0.49\textwidth}
  \centering
  \includegraphics[width=\linewidth]{unc_vs_error.png}
  \caption{Case-level analysis. Left: higher mean uncertainty correlates
  with lower Dice. Right: per-image Unc-AUROC remains high even for
  difficult cases---uncertainty is most informative where it matters most.}
  \label{fig:unc_vs_error}
\end{subfigure}
\caption{Deferral behavior and case-level uncertainty analysis. Left panel
shows TTA deferral modes; right panel confirms that uncertainty quality
holds across the difficulty spectrum.}
\label{fig:deferral_views}
\end{figure}

\subsection{Risk-coverage analysis}
\label{sec:results_risk_coverage}

\textbf{Claim:} Selective prediction using TTA uncertainty produces a
steeper risk-coverage curve than MC Dropout and achieves clinically
acceptable Dice ($\geq 0.82$) at 90\% coverage.

Figure~\ref{fig:risk_coverage} shows Dice as a function of coverage for both
uncertainty methods. At full coverage (no deferral), TTA Dice is 0.768 and MC
Dropout is 0.764. As coverage decreases (more pixels deferred), both improve,
but TTA improves faster.

At 90\% coverage (deferring 10\% of pixels):
\begin{itemize}[leftmargin=*]
  \item MC Dropout: Dice = 0.778
  \item TTA: Dice = 0.849
\end{itemize}

The clinical threshold for retinal vessel segmentation is typically taken as
Dice $\geq 0.82$ \citep{fraz2012ensemble}. TTA reaches this threshold at
roughly 92.7\% coverage, while MC Dropout does not reach it until roughly
71.8\% coverage. This gap is practically significant: the TTA operating point
retains most pixels while clearing the clinical bar, whereas MC Dropout
requires much heavier abstention.

The AUCC values quantify this difference:
\begin{itemize}[leftmargin=*]
  \item TTA AUCC (Dice): 0.846 vs.\ MC Dropout: 0.801
  \item TTA AUCC (AUC): 0.888 vs.\ MC Dropout: 0.882
\end{itemize}

\begin{figure}[H]
\centering
\includegraphics[width=\textwidth]{risk_coverage.png}
\caption{Risk-coverage curves for MC Dropout and TTA. TTA reaches
clinically acceptable Dice (0.82) at substantially higher coverage than
MC Dropout, requiring far less human review to meet the clinical bar.}
\label{fig:risk_coverage}
\end{figure}

\paragraph{Operating point analysis.} Three clinically motivated operating
points illustrate the practical gap. At the \emph{clinical target}
(Dice~$\geq 0.82$), TTA retains 90\% of pixels while MC Dropout must
drop to 80\%~coverage---a 10-point gap at the same quality bar.
At the balanced-F1 operating point, MC Dropout reaches Dice~0.855 but
only at 71.9\% coverage; TTA achieves Dice~0.838 at 80\% coverage.
Under aggressive deferral (75\% coverage), TTA pushes Dice to 0.867.
The consistent pattern: TTA reaches every quality target at
higher coverage, meaning less human review for the same clinical
guarantee.

\subsection{Calibration and deferral quality are decoupled}
\label{sec:results_calibration}

\textbf{Claim:} Better calibration does not produce better deferral.
Temperature scaling worsens the reported ECE values and leaves deferral
operating points unchanged, demonstrating that calibration metrics and
selective prediction quality measure different properties.

\emph{Insight before the numbers.} Temperature scaling is a monotonic
rescaling: it cannot change the order of pixels by uncertainty. Any
decision rule that thresholds on that order is therefore invariant to
temperature scaling by construction. Table~\ref{tab:calibration} confirms
this prediction empirically.

\begin{table}[H]
\centering
\caption{Effect of temperature scaling on calibration and deferral. Both ECE
values \emph{increase} after calibration while Unc-AUROC is unchanged,
confirming that rank-preserving transforms cannot improve a discriminative
metric. $T$ is the learned temperature.}
\label{tab:calibration}
\scriptsize
\setlength{\tabcolsep}{4pt}
\renewcommand{\arraystretch}{1.05}
\begin{center}
\begin{tabular}{l c c c c c}
\toprule
\textbf{Method} & $T$ & \textbf{ECE (pre)} & \textbf{ECE (post)} &
$\Delta$\textbf{ECE} & \textbf{Unc-AUROC} \\
\midrule
MC Dropout & 1.20 & \textbf{0.035} & 0.043 & $+0.008$ & 0.722 \\
TTA        & 1.35 & \textbf{0.035} & 0.039 & $+0.005$ & \textbf{0.881} \\
\bottomrule
\end{tabular}
\end{center}
\end{table}

For MC Dropout, temperature scaling increases ECE from 0.035 to 0.043
while leaving uncertainty ranking unchanged at 0.722. For TTA, ECE rises
from 0.035 to 0.039; Unc-AUROC remains at 0.881. Neither temperature-scaled
configuration improves downstream error reduction.

\paragraph{Why this happens.} Both models learn $T > 1$, confirming mild
overconfidence. But ECE is already below 0.04 before calibration, leaving
little room for improvement. The learned temperatures rescale probabilities
without changing uncertainty \emph{rankings}, so discriminative metrics
stay flat while calibration values shift unfavorably. The lesson is both
concrete---evaluate temperature scaling by Unc-AUROC and operating points,
not ECE alone---and conceptual: calibration optimizes a distributional
property; deferral optimizes a discriminative one. A monotonic rescaling
cannot move the latter, which is exactly why the two objectives decouple.

\begin{figure}[H]
\centering
\includegraphics[width=0.49\textwidth]{reliability_mc.png}
\includegraphics[width=0.49\textwidth]{reliability_tta.png}
\caption{Reliability diagrams for MC Dropout (left) and TTA (right), with
and without temperature scaling. Both models are already close to the
diagonal before calibration, leaving little room for temperature scaling to
improve.}
\label{fig:reliability}
\end{figure}

\subsection{Cross-dataset evaluation}
\label{sec:results_cross_dataset}

\textbf{Claim:} Uncertainty quality transfers more robustly than segmentation
accuracy under domain shift.

\emph{Insight before the numbers.} If errors on unfamiliar data concentrate
in the same regions the model was already uncertain about, the
uncertainty--error correlation survives the shift even when the model's
raw accuracy does not. Table~\ref{tab:cross_dataset} shows exactly that
pattern on the zero-shot DRIVE$\to$STARE/CHASE\_DB1 transfer.

\begin{table}[H]
\centering
\caption{Cross-dataset evaluation (zero-shot transfer from DRIVE).
Segmentation quality drops by 6--7 Dice points under domain shift, yet
uncertainty quality (Unc-AUROC) holds within 1--2 points of the in-domain
value. Uncertainty transfers more robustly than segmentation---deferral
remains reliable exactly when it matters most.}
\label{tab:cross_dataset}
\scriptsize
\setlength{\tabcolsep}{3.5pt}
\renewcommand{\arraystretch}{1.03}
\begin{center}
\begin{tabular}{l c c c c c}
\toprule
\textbf{Dataset} & \textbf{Dice} & \textbf{AUC} & \textbf{ECE} &
\textbf{Unc-AUROC} & \textbf{AUCC (Dice)} \\
\midrule
DRIVE (in-domain) & \textbf{0.791} & \textbf{0.960} & 0.043 & \textbf{0.850} & \textbf{0.793} \\
STARE             & 0.727 & 0.949 & 0.038 & 0.846 & 0.773 \\
CHASE\_DB1        & 0.726 & 0.952 & \textbf{0.035} & 0.835 & 0.766 \\
\bottomrule
\end{tabular}
\end{center}
\end{table}

Dice drops from 0.791 (DRIVE) to 0.727 (STARE) and 0.726 (CHASE\_DB1)---a
roughly 6.5-point decline. Unc-AUROC, however, remains strong: 0.846 on
STARE and 0.835 on CHASE\_DB1, both close to the in-domain value. This
result is explained by the fact that on out-of-distribution data, the model
makes more errors in regions where it is also more uncertain (unfamiliar
appearance patterns, different vessel calibers), preserving the
uncertainty--error correlation.

This is the clearest clinical argument in the paper: on unfamiliar data
where deferral matters most, the uncertainty signal is at its most
reliable. A system deployed on a new scanner or patient population can
lean on the decision layer exactly when the estimation layer is most
stressed. Robustness, for a deferral-enabled pipeline, is a property of
the \emph{joint system}, not of the segmentation head alone.

\begin{figure}[H]
\centering
\includegraphics[width=\linewidth,height=0.4\textheight,keepaspectratio]{qualitative_panel.png}
\caption{Qualitative results across datasets. Errors concentrate at vessel
boundaries (column d). TTA uncertainty tracks these regions (column e), and
the deferral policy routes them for review (column f) while retaining the
large majority of correctly segmented pixels.}
\label{fig:qualitative}
\end{figure}

\subsection{Distribution shift robustness}
\label{sec:results_shift}

We also evaluate robustness to synthetic distribution shifts: Gaussian noise
and blur at two severity levels.

\begin{table}[H]
\centering
\caption{Effect of synthetic distribution shift. Uncertainty degrades more
slowly than segmentation: Dice falls from 0.764 to 0.186 under strong shift
($-$75\% relative) while Unc-AUROC only drops from 0.765 to 0.655 ($-$14\%
relative). The clean baseline is in \textbf{bold}; the strong-shift row is in
\emph{italics}.}
\label{tab:shift}
\scriptsize
\setlength{\tabcolsep}{3.5pt}
\renewcommand{\arraystretch}{1.03}
\begin{center}
\begin{tabular}{l c c c c c}
\toprule
\textbf{Shift level} & \textbf{Dice} & \textbf{AUC} & \textbf{ECE} &
\textbf{Unc-AUROC} & \textbf{ERR (adaptive)} \\
\midrule
None              & \textbf{0.764} & \textbf{0.965} & \textbf{0.037} & \textbf{0.765} & \textbf{45.6\%} \\
Mild              & 0.517 & 0.840 & 0.056 & 0.648 & 7.9\% \\
\emph{Strong}     & \emph{0.186} & \emph{0.758} & \emph{0.094} & \emph{0.655} & \emph{25.0\%} \\
\bottomrule
\end{tabular}
\end{center}
\end{table}

Under mild shift ($\sigma_{\text{noise}} = 0.02$, $3\times3$ Gaussian blur
with $\sigma = 0.5$), Dice drops to 0.517---a 32\% relative decline---while
Unc-AUROC drops to 0.648, still above chance but substantially degraded.
Under strong shift ($\sigma_{\text{noise}} = 0.10$, $7\times7$ blur,
4$\times$ resolution reduction), Dice collapses to 0.186 and Unc-AUROC is
0.655.

The asymmetry is notable: even under catastrophic model failure (Dice 0.186),
uncertainty maintains above-random discriminative power (AUROC 0.655 vs.\
0.5 random). However, deferral error reduction under strong shift is limited
because most pixels are wrong and the uncertainty signal, while above random,
is not sharp enough to efficiently separate remaining correct predictions.

This motivates a concrete deployment rule: monitor both average uncertainty
magnitude and Unc-AUROC on a small labeled subset. When Unc-AUROC drops
below $\sim$0.70, switch from pixel-level deferral to image-level flagging
---the uncertainty signal is too noisy for fine-grained triage but still
informative enough to rank whole images by expected error.
